\documentclass[11pt]{article}

\usepackage[margin=1in]{geometry}
\usepackage{amsmath,amssymb,amsthm}
\usepackage{graphicx}
\usepackage[hidelinks]{hyperref}

\newtheorem{theorem}{Theorem}
\newtheorem{corollary}{Corollary}
\newtheorem{remark}{Remark}

\newcommand{\CWOS}{\textit{CWO-S}}
\newcommand{\eps}{\varepsilon}

\title{An Arbitrary \(\ell_1\)-Sum Extension for Property \(\CWOS\)}
\author{Run fa\_banach\_001, lane 7}
\date{2026-06-20}

\begin{document}
\maketitle

\section*{Status}

This packet gives a likely valid partial result for the natural question of
Haller--Kuuseok--Poldvere~\cite{HKP25}.  The source paper proves that
\(L_1(\mu,X)\) has property \(\CWOS\) whenever \(\mu\) is non-zero
\(\sigma\)-finite and \(X\) is weakly uniformly rotund, and asks whether
``weakly uniformly rotund'' can be replaced by ``has property \(\CWOS\)''.

The result here proves the positive answer on the purely atomic side in a
slightly stronger form: arbitrary, not merely countable, \(\ell_1\)-sums of
\(\CWOS\) spaces have \(\CWOS\).  Consequently \(L_1(\mu,X)\) has \(\CWOS\)
for arbitrary purely atomic measure spaces whenever \(X\) has \(\CWOS\).  The
non-atomic case, and the injective tensor-product question, remain open here.

\section*{Source Problem Evidence}

The source question appears in the introduction of~\cite{HKP25}, page 3 of the
arXiv PDF:

\begin{center}
\includegraphics[width=0.95\linewidth]{figures/open_problem_crop.png}
\end{center}

\section*{Transcription}

After proving the \(w\)UR-valued \(L_1\)-theorem and the countable
\(\ell_1\)-sum stability proposition, the authors ask whether the theorem
remains true when \(X\) is assumed to be any Banach space with property
\(\CWOS\), and moreover whether the injective tensor product of two Banach
spaces with property \(\CWOS\) must have property \(\CWOS\).

\section*{Definitions}

A Banach space \(X\) has property \(\CWOS\) if, for every finite convex
combination
\[
  C=\sum_{k=1}^n \lambda_k U_k
\]
of relatively weakly open subsets \(U_k\) of \(B_X\), every point of
\(C\cap S_X\) is an interior point of \(C\) in the relative weak topology of
\(B_X\).

For a family \((X_i)_{i\in I}\), write
\[
  \left(\bigoplus_{i\in I} X_i\right)_1
\]
for the Banach space of all families \(x=(x_i)\) with
\(\sum_{i\in I}\|x_i\|<\infty\), normed by that sum.

\section*{Proof Intuition}

The source paper proves countable \(\ell_1\)-sum stability.  In an arbitrary
\(\ell_1\)-sum, any two fixed vectors \(x,y\) have countable support, so their
norm-one convex combination \(z=\lambda x+\nu y\) lives in a countable
one-complemented summand \(A\).  That summand has \(\CWOS\) by the source
proposition.  A weakly small perturbation \(w\) of \(z\) in the whole
arbitrary sum may have a small component \(b\) outside \(A\).  Decompose the
normalized \(A\)-part of \(w\) inside \(A\), then put the same off-support
component \(b\) into both summands.  The \(\ell_1\)-norm budget absorbs this
extra component exactly because \(\|P_Aw\|+\|(I-P_A)w\|=\|w\|\le1\).

\begin{theorem}\label{thm:arbitrary-l1}
Let \((X_i)_{i\in I}\) be an arbitrary family of Banach spaces with property
\(\CWOS\).  Then
\[
  Z=\left(\bigoplus_{i\in I} X_i\right)_1
\]
has property \(\CWOS\).
\end{theorem}

\begin{proof}
We use the standard reduction to convex combinations of two relatively weakly
open sets, as in Lemma 4.1(b) of~\cite{ABHLP} and in the proof of Proposition
1.2 of~\cite{HKP25}.  Thus let \(U,V\) be relatively weakly open subsets of
\(B_Z\), let \(\lambda,\nu>0\) with \(\lambda+\nu=1\), and suppose
\[
  x\in U,\qquad y\in V,\qquad z=\lambda x+\nu y\in S_Z.
\]
It is enough to find a relative weak neighbourhood \(W\) of \(z\) in \(B_Z\)
such that \(W\subset \lambda U+\nu V\).

Choose a finite set \(F\subset S_{Z^*}\) and \(\eta>0\) such that
\[
  U_0=\{u\in B_Z: |f(u-x)|<3\eta \text{ for all } f\in F\}\subset U
\]
and similarly
\[
  V_0=\{v\in B_Z: |f(v-y)|<3\eta \text{ for all } f\in F\}\subset V.
\]
Let \(J=\operatorname{supp}x\cup\operatorname{supp}y\).  This set is
countable.  Put
\[
  A=\left(\bigoplus_{i\in J}X_i\right)_1,\qquad
  B=\left(\bigoplus_{i\in I\setminus J}X_i\right)_1,
\]
so that \(Z=A\oplus_1 B\), and let \(P:Z\to A\) be the norm-one coordinate
projection.  The vectors \(x,y,z\) belong to \(A\).  By Proposition 1.2 of
\cite{HKP25}, \(A\) has property \(\CWOS\).

For \(f\in F\), denote by \(f_A\) and \(f_B\) its restrictions to \(A\) and
\(B\).  Define relatively weakly open neighbourhoods of \(x\) and \(y\) in
\(B_A\) by
\[
  U_A=\{a\in B_A: |f_A(a-x)|<\eta \text{ for all } f\in F\},
\]
\[
  V_A=\{a\in B_A: |f_A(a-y)|<\eta \text{ for all } f\in F\}.
\]
Since \(A\) has \(\CWOS\) and \(z=\lambda x+\nu y\in S_A\), there is a
relative weak neighbourhood \(W_A\) of \(z\) in \(B_A\) such that
\[
  W_A\subset \lambda U_A+\nu V_A.
\]
Choose a finite set \(G\subset A^*\) and \(\delta>0\) such that
\[
  \{a\in B_A: |g(a-z)|<\delta \text{ for all } g\in G\}\subset W_A.
\]
Let \(h\in S_{A^*}\) satisfy \(h(z)=1\).  Decreasing \(\rho>0\) if necessary,
assume
\[
  \rho<\eta,\qquad
  \rho<\delta/4,\qquad
  \rho(1+\|g\|)<\delta/2 \quad(g\in G).
\]
Now set
\[
  W=\left\{w\in B_Z:
  |g(Pw-z)|<\rho \ (g\in G),\quad |h(Pw-z)|<\rho\right\}.
\]
This is a relative weak neighbourhood of \(z\) in \(B_Z\).

Fix \(w\in W\).  Write \(a=Pw\in A\) and \(b=w-a\in B\), and set
\(t=\|a\|\).  Since \(h(a)>1-\rho\), we have \(t>1-\rho>0\), and
\(\|b\|\le 1-t<\rho\).  For \(g\in G\),
\[
  \left|g\left(\frac{a}{t}-z\right)\right|
  \le \frac{|g(a-z)|}{t}+\left|\frac1t-1\right|\,|g(z)|
  <2\rho+2\rho\|g\|<\delta.
\]
Thus \(a/t\in W_A\).  Hence there exist \(a_1\in U_A\) and \(a_2\in V_A\)
such that
\[
  \frac{a}{t}=\lambda a_1+\nu a_2.
\]
Define
\[
  u=ta_1+b,\qquad v=ta_2+b
\]
as elements of \(A\oplus_1B=Z\).  Then
\[
  \lambda u+\nu v=t(\lambda a_1+\nu a_2)+b=a+b=w.
\]
Also
\[
  \|u\|\le t\|a_1\|+\|b\|\le t+\|b\|=\|w\|\le1,
\]
and similarly \(\|v\|\le1\).

It remains to check \(u\in U_0\) and \(v\in V_0\).  For \(f\in F\),
\[
  f(u-x)=t f_A(a_1-x)+(t-1)f_A(x)+f_B(b).
\]
Since \(a_1\in U_A\), \(\|f_A\|\le1\), \(\|x\|\le1\), \(|1-t|<\rho\), and
\(\|b\|<\rho\), we get
\[
  |f(u-x)|<\eta+\rho+\rho<3\eta.
\]
Thus \(u\in U_0\subset U\).  The same estimate with \(a_2,y,V_A\) gives
\(v\in V_0\subset V\).  Therefore \(w\in\lambda U+\nu V\).  Since
\(w\in W\) was arbitrary, \(W\subset\lambda U+\nu V\), completing the proof.
\end{proof}

\begin{corollary}\label{cor:atomic}
Let \(\mu\) be a purely atomic measure whose finite-measure atoms generate the
\(L_1\)-space, and let \(X\) be a Banach space with property \(\CWOS\).  Then
\(L_1(\mu,X)\) has property \(\CWOS\).
\end{corollary}

\begin{proof}
On each finite-measure atom \(A\), every Bochner integrable function is
almost everywhere constant.  Therefore \(L_1(\mu,X)\) is isometric to the
\(\ell_1\)-sum over the finite-measure atoms, with each summand isometric to
\(X\) after the harmless scaling by \(\mu(A)\).  Apply
Theorem~\ref{thm:arbitrary-l1}.
\end{proof}

\section*{Verification Notes}

No computation is needed.  The proof uses only:
\begin{itemize}
  \item the source paper's countable \(\ell_1\)-sum stability theorem;
  \item the standard two-set reduction for property \(\CWOS\);
  \item the elementary fact that every element of an arbitrary \(\ell_1\)-sum
        has countable support;
  \item the norm identity in \(A\oplus_1 B\).
\end{itemize}

\section*{Limitations}

This is a partial result.  For the source paper's \(\sigma\)-finite setting,
the purely atomic case is already morally countable and close to Proposition
1.2 of~\cite{HKP25}.  The new contribution is the arbitrary-index
\(\ell_1\)-sum argument and the clean atomic formulation without a countability
restriction.  The non-atomic \(L_1(\mu,X)\) case for general \(\CWOS\) spaces
and the injective tensor-product question are not settled by this packet.

\section*{Novelty and Literature Check}

I searched the run indexes for arXiv:2503.09182, \(\CWOS\), \(\ell_1\)-sum,
injective tensor product, and convex combinations of relatively weakly open
sets.  The only overlapping run packet was a prior partial result for the
older CWO tensor question from arXiv:1806.10693; it does not contain the
arbitrary \(\ell_1\)-sum upgrade proved here.

I also ran bounded web searches for ``arbitrary 1-sum property CWO-S'',
``l1-sum CWO-S Banach spaces'', and related direct-sum phrases.  The searches
found the source paper's countable result and adjacent CWO literature, but no
exact arbitrary-index statement.  Novelty confidence is moderate: the proof is
short and natural from the source proposition, so it may be folklore.

\section*{Human Review Recommendation}

Check the passage from the countable support summand \(A\) back to the full
sum \(A\oplus_1B\), especially the normalization \(a/\|a\|\in W_A\) and the
estimate showing \(u=ta_1+b\) and \(v=ta_2+b\) stay in the original weak
neighbourhoods.  This is the only genuinely new step beyond the source
paper's countable theorem.

\begin{thebibliography}{9}

\bibitem{HKP25}
R. Haller, P. Kuuseok, and M. Poldvere,
\emph{Norm-one points in convex combinations of relatively weakly open subsets
of the unit ball in the spaces \(L_1(\mu,X)\)},
arXiv:2503.09182, 2025.
\url{https://arxiv.org/abs/2503.09182}

\bibitem{ABHLP}
T. A. Abrahamsen, J. Becerra Guerrero, R. Haller, V. Lima, and M. Poldvere,
\emph{Banach spaces where convex combinations of relatively weakly open
subsets of the unit ball are relatively weakly open},
Studia Math. \textbf{250} (2020), no. 3, 297--320.
\url{https://arxiv.org/abs/1806.10693}

\end{thebibliography}

\end{document}
